\chapter{CONCLUSIONS ET PERSPECTIVES}\label{chap:conclusions-perspectives}
\thispagestyle{empty}

\section{Bilan general du projet}

Ce projet avait pour objectif de concevoir, evaluer et integrer un systeme de resume automatique d'actualites multilingues. Le travail realise ne s'est pas limite a l'entrainement d'un modele de generation de texte. Il a reuni plusieurs dimensions complementaires: l'etude des modeles de resume, l'entrainement de modeles mBART, l'utilisation d'un modele mT5 pret a l'emploi comme reference externe, la construction d'une evaluation multidimensionnelle et le developpement d'une application web capable de collecter, traiter, resumer et afficher des actualites provenant de flux RSS reels.

La problematique initiale etait de savoir comment evaluer des modeles de resume multilingue selon des criteres plus larges que la simple similarite lexicale avec une reference, puis comment utiliser ces modeles dans un systeme applicatif. Les chapitres precedents ont montre que cette question demande une approche de bout en bout. Un bon modele de resume doit produire un texte court et lisible, mais il doit aussi rester fidele a la source, eviter les repetitions, couvrir les informations importantes et fonctionner avec une latence compatible avec un flux d'actualites.

Le systeme final repond a cette problematique par une architecture en plusieurs couches. Les donnees academiques, principalement XL-Sum, ont ete utilisees pour entrainer et evaluer les modeles. Les modeles mBART entraines dans le cadre du projet ont ensuite ete compares avec le modele \texttt{mt5-xlsum}. Enfin, l'application Flask a permis d'integrer ces modeles dans un scenario pratique, avec collecte RSS, extraction du corps des articles, nettoyage, generation de resumes, stockage SQLite, ingestion en arriere-plan et interface filtrable.

\section{Synthese des resultats obtenus}

L'evaluation complete a ete realisee sur le split \texttt{test} de \texttt{csebuetnlp/xlsum}, avec 52.219 exemples repartis sur 11 langues executables. Les langues allemand, italien, neerlandais et roumain etaient prevues dans le plan applicatif, mais elles n'ont pas ete evaluees quantitativement car les sous-ensembles publics correspondants n'etaient pas disponibles dans la version locale de XL-Sum. Cette distinction est importante: l'application conserve un support multilingue plus large, mais les conclusions numeriques solides concernent les 11 langues effectivement evaluees.

Les resultats montrent que les trois modeles compares sont proches en score global de capacite. Le modele \texttt{mt5-xlsum} obtient la meilleure moyenne globale, avec un score \texttt{capability\_overall} de 0.6631. Le modele \texttt{mbart50\_xlsum}, entraine dans le cadre du projet, atteint 0.6578, et \texttt{mbart-xlsum-2} atteint 0.6570. L'ecart entre le meilleur modele et le meilleur modele entraine dans le projet reste donc faible, ce qui montre que les modeles mBART developpes sont competitifs face a une reference mT5 deja specialisee.

La comparaison montre cependant que les modeles n'ont pas le meme profil. \texttt{mt5-xlsum} est le plus fort pour les metriques classiques de similarite avec les resumes humains, notamment ROUGE-L, BLEU et METEOR-lite. Ce resultat est coherent avec son role de modele de comparaison deja entraine pour le resume multilingue. A l'inverse, les deux modeles mBART offrent un avantage applicatif clair sur la latence et l'efficacite. \texttt{mbart50\_xlsum} represente le compromis le plus prudent entre capacite globale, vitesse et fidelite a la source. \texttt{mbart-xlsum-2} est le plus rapide en moyenne et obtient une meilleure completude automatique, mais il est aussi plus abstractive et moins directement ancre dans le texte source.

Ces resultats confirment que le choix du modele ne doit pas etre reduit a un classement unique. Si l'objectif principal est la proximite avec les resumes de reference, \texttt{mt5-xlsum} constitue le meilleur choix. Si l'objectif est une integration applicative stable et plus rapide, \texttt{mbart50\_xlsum} est un candidat solide comme modele par defaut. Si la priorite est la vitesse et une reformulation plus forte, \texttt{mbart-xlsum-2} peut etre utilise comme alternative, avec une attention particuliere aux articles sensibles du point de vue factuel.

\section{Apports du travail}

Le premier apport du projet est la mise en place d'une chaine complete pour le resume d'actualites multilingues. Le travail couvre l'analyse du probleme, la revue des approches existantes, l'entrainement de modeles mBART, l'evaluation experimentale et l'integration dans une application web. Cette continuite entre recherche experimentale et implementation logicielle donne au projet une valeur pratique: les modeles ne sont pas seulement compares sur un corpus fixe, ils sont aussi utilises dans un systeme capable de traiter des sources reelles.

Le deuxieme apport concerne l'evaluation. Les metriques classiques comme ROUGE, BLEU et METEOR-lite ont ete completees par des metriques de comportement et des scores proxy composites. Les dimensions de coherence, exactitude, clarte, pertinence, efficacite, fidelite a la source et completude permettent d'analyser plus finement les sorties des modeles. Cette approche ne remplace pas une evaluation humaine, mais elle evite de tirer des conclusions uniquement a partir d'un score lexical.

Le troisieme apport est architectural. L'application separe la generation des resumes du moment de la requete utilisateur. Les articles sont collectes et resumes en arriere-plan, puis les resultats sont stockes dans SQLite. Cette decision rend le systeme plus utilisable, car l'utilisateur consulte des resumes deja prepares au lieu d'attendre l'execution d'un modele a chaque appel API\@. La presence de filtres par langue, sujet, pays, region, source, mot-cle et modele rend aussi le prototype plus proche d'un outil exploitable.

Enfin, le projet prend en compte des contraintes pratiques et ethiques. Les liens vers les articles originaux sont conserves, les modeles multilingues sont controles selon la langue choisie, et des mecanismes de nettoyage, de post-traitement et de resume extractif de secours limitent les sorties de mauvaise qualite. Le resume est donc presente comme une aide a la lecture rapide, non comme un remplacement definitif de l'article source.

\section{Limites du projet}

Malgre ces resultats, plusieurs limites doivent etre soulignees. La premiere concerne l'evaluation automatique. Les scores proxy construits dans le projet sont utiles pour comparer les modeles dans les memes conditions, mais ils restent heuristiques. Ils ne prouvent pas a eux seuls la qualite humaine d'un resume, ni sa veracite factuelle dans le monde reel. Une evaluation humaine reste necessaire pour mesurer plus directement la coherence, la fluidite, la pertinence et la confiance percue par les utilisateurs.

La deuxieme limite concerne la fidelite factuelle. Le projet mesure la proximite au texte source avec des signaux comme \texttt{source\_coverage}, \texttt{fragment\_coverage} et \texttt{quality\_factuality}. Ces mesures peuvent indiquer si le resume reste proche de l'article, mais elles ne constituent pas un systeme complet de verification des faits. Un resume peut etre lexicalement proche de la source tout en supprimant un contexte important, ou reformuler une information de maniere ambigue.

La troisieme limite vient de la difference entre XL-Sum et les flux RSS reels. XL-Sum fournit des articles et des resumes de reference propres, alors que les sources d'actualites en ligne contiennent souvent du bruit, des structures HTML variables, des titres repetes, des extraits incomplets ou des contenus non journalistiques. L'application integre des mecanismes de nettoyage et de filtrage, mais la qualite de l'extraction depend encore fortement des sites sources.

La quatrieme limite est linguistique. L'evaluation finale couvre 11 langues, alors que l'application en propose 15. Les langues allemand, italien, neerlandais et roumain restent disponibles dans le flux applicatif, mais elles ne disposent pas ici de resultats experimentaux comparables. De plus, les performances varient selon les langues, notamment en fonction de la taille des donnees disponibles, de la tokenisation et des specificites de chaque systeme d'ecriture.

Enfin, la latence mesuree depend de l'environnement materiel, de la taille des batchs, des reglages de generation et de l'etat de la machine au moment de l'evaluation. Les conclusions sur la vitesse sont donc valables pour la configuration experimentale utilisee, mais elles doivent etre revalidees avant un deploiement sur une infrastructure differente.

\section{Perspectives d'amelioration}

Une premiere perspective consiste a ajouter une evaluation humaine structuree. Des annotateurs pourraient comparer les resumes selon plusieurs criteres: fidelite a la source, couverture de l'information, clarte, absence de repetition et utilite pour une lecture rapide. Cette evaluation permettrait de verifier si les scores proxy sont bien alignes avec le jugement humain et d'ajuster leurs poids si necessaire.

Une deuxieme perspective concerne l'evaluation factuelle. Le systeme pourrait integrer des metriques plus specialisees de coherence factuelle ou des methodes basees sur questions-reponses et inference textuelle. L'objectif serait de detecter plus directement les contradictions, les entites incorrectes, les chiffres modifies ou les informations absentes de la source. Cette evolution serait particulierement importante pour les nouvelles politiques, economiques, medicales ou scientifiques.

Une troisieme piste est l'amelioration des modeles. Les deux modeles mBART pourraient etre reentraines avec une validation plus fine par langue, des jeux de donnees supplementaires, des strategies d'adaptation legere ou une distillation vers des modeles plus petits. Une approche interessante serait de conserver la qualite du modele le plus fort tout en reduisant le cout de generation, afin de mieux servir l'application en temps reel.

Une quatrieme perspective est la construction d'un jeu de test local pour les langues non couvertes dans l'evaluation XL-Sum finale. Un ensemble d'articles et de resumes de reference pour l'allemand, l'italien, le neerlandais et le roumain permettrait de valider quantitativement les 15 langues de l'application. Ce jeu de test pourrait aussi inclure des exemples directement issus des sources RSS utilisees par le systeme, afin de rapprocher l'evaluation du contexte applicatif reel.

Une cinquieme perspective concerne l'application web. Le systeme pourrait proposer une comparaison visuelle entre plusieurs resumes du meme article, un indicateur de confiance, une alerte lorsque l'extraction du texte est faible, ou une option permettant a l'utilisateur de signaler un resume incorrect. Les retours utilisateurs pourraient ensuite alimenter une evaluation continue et aider a identifier les sources ou les langues qui posent le plus de problemes.

Enfin, l'architecture pourrait etre renforcee pour un usage plus large. L'utilisation d'une base plus robuste que SQLite, l'ajout d'une file de taches dediee, la mise en cache plus fine des modeles, la surveillance des erreurs d'ingestion et la gestion de quotas par source permettraient de rendre le systeme plus fiable dans un contexte de production.

\section{Conclusion finale}

En conclusion, ce projet montre qu'un systeme de resume automatique d'actualites multilingues doit etre pense comme un ensemble complet, et non comme un simple appel a un modele de langage. La qualite finale depend a la fois des donnees, de l'entrainement, des metriques d'evaluation, de la fidelite au texte source, du temps de generation et de l'integration logicielle.

Les resultats obtenus montrent que le modele mT5 pret a l'emploi reste une reference forte pour la similarite avec les resumes humains, mais que les modeles mBART entraines dans le cadre du projet offrent un compromis applicatif pertinent, en particulier pour la vitesse et l'utilisation dans un flux d'actualites. Le modele \texttt{mbart50\_xlsum} apparait comme le choix le plus prudent pour une utilisation par defaut, tandis que \texttt{mbart-xlsum-2} constitue une alternative plus rapide et plus abstractive.

Le projet atteint ainsi son objectif principal: proposer une chaine de bout en bout qui relie l'entrainement et l'evaluation de modeles de resume multilingue a une application web fonctionnelle. Les limites identifiees ouvrent des pistes claires pour la suite, notamment l'evaluation humaine, le controle factuel, l'extension linguistique et le passage vers une architecture plus robuste. Ces perspectives montrent que le systeme developpe constitue une base exploitable pour de futurs travaux sur l'acces rapide, multilingue et fiable a l'information d'actualite.
